\chapter{共价键：共享电子不是平均分家}

\begin{kaichang}
医院里有一种蓝色的钢瓶，装着压缩氧气。登山者背它上珠峰，潜水员背的则是氧气和别的气体的混合瓶。此刻你吸进的这一口空气里，大约五分之一是氧气——它们以“两个氧原子抱成一团”的形式飘浮着，从肺泡钻进你的血液。

再看一眼你手边的另一些东西：塑料笔杆、矿泉水、橡皮筋、你自己的头发。它们和氧气、和食盐都不一样——食盐一摔就碎、丢进水里就消失、熔化了能导电（第 18 章）；而塑料拉一拉会先变形，水在锅里烧开变成气也仍旧是水，头发烧起来有股焦味而不是“熔化”。

现在猜一猜：一个氧原子为什么非要和另一个氧原子粘在一起？两个一模一样的原子，谁的“抢电子”力气都一样大——食盐那种“一个交出、一个收下”的办法（第 18 章）在这里根本走不通。它们到底是怎样结伴的？这一章结束时，我们要用新模型把这只蓝色钢瓶、这杯水、这根塑料尺重新解释一遍。
\end{kaichang}

\section{当“交出电子”这条路走不通}

先回忆一下第 18 章的离子键是怎样成立的：钠最外层只有 1 个电子，丢起来“便宜”；氯最外层 7 个电子，抢到 1 个就凑满外层。一丢一抢，两边的外层排布都变得更省心，整个体系的总能量下降，于是电子完成转移，正负离子靠库仑吸引堆成晶体。这笔账能算平，前提是\textbf{双方对电子的胃口差得足够多}。

可是生活里有大量物质不是这样组成的。氢气和氢气、氧气和氧气、氯气和氯气——两个完全相同的原子相遇，谁也没法从对方手里白拿电子：你抢我的，我也以同样的力气抢你的，僵持不下。碳和氢、碳和氧、氮和氢之间也类似：它们对电子的胃口有差别，但差别不够大，谁也做不到把对方的电子整个没收。

然而这些物质又确确实实稳定存在：两个氢原子会结结实实地粘成氢分子 \chem{H_2}，粘得相当牢——想把它们拆开，需要相当可观的能量。宏观上的证据到处都是：氢气能在空气中安静地存放（只要别点火），氧气在钢瓶里一存几年，塑料袋提着重物不撒。原子之间显然存在一种和第 18 章完全不同的“结伴方式”。

把能观察到的事实摆在一起对比，差别一目了然：食盐熔点超过 $800\,^{\circ}\mathrm{C}$，固态不导电，熔融或溶解后导电；而糖（由碳、氢、氧组成）加热到 $190\,^{\circ}\mathrm{C}$ 左右就开始变色分解，根本不给你熔化的机会，它的水溶液也几乎不导电。蜡烛一烤就软，冰一加热就成水。这些物质的共同点是：它们的“基本队伍”是一个个独立的小集团——\textbf{分子}，集团内部原子粘得很牢，集团与集团之间却拉得很松（后者是第 22 章的话题）。第 18 章那种“电子转移加静电堆积”的模型，解释不了分子内部这股牢劲。

\section{共享：两个原子核同时抓住一对电子}

把镜头推进到两个正在靠近的氢原子。每个氢原子只有一个质子和一个电子。第 17 章已经搭好了舞台：两个原子靠近时，核与核互相排斥、电子与电子互相排斥，但每个电子同时也被对方的原子核吸引。体系的总能量随距离变化，画出来是一条先下降后上升的势能曲线——曲线有个“谷底”，谷底对应的距离就是键长，谷底的深度就是键能。

现在问关键问题：对两个氢原子来说，谷底为什么会存在？为什么能量不是一路排斥、越近越高？

答案藏在电子的“双重身份”里。当两个氢原子靠得足够近，两个电子的活动范围开始重叠。此时\textbf{每个电子都能同时感受到两个原子核的吸引}。用第 9 章的电子云语言说：两团电子云叠在了一起，两个核之间的区域里电子出现的概率明显增大，电子不再是“谁的私有财产”，而是被两个核共同“租用”。一个电子被两个正电荷同时吸引，比只被一个核吸引时能量更低；两股吸引合在一起，足以盖过核—核、电子—电子的排斥——于是势能曲线上出现了谷底。这一对被共同占用的电子，就是\textbf{共价键}。

请特别注意这个图景和第 18 章的区别。离子键里，电子完成了“过户”：钠的电子搬到氯那边，之后各过各的，靠静电相吸。共价键里，电子没有过户，而是“共有”：两个原子核对这对电子的吸引把它们拴在了一起，谁也别想独占。可以打一个不太严谨但好记的比方：离子键像一方把房子卖了拿钱走人，共价键像两家合买一套房，谁搬走房子都不再成立。

同样的逻辑立刻解释了一大批物质。两个氯原子各出 1 个电子共享，各自外层“看起来”都凑满了 8 个，形成 \chem{Cl_2}；一个氧原子和两个氢原子分别共享电子对，形成水分子 \chem{H_2O}；碳原子最外层 4 个电子，可以和四个氢各共享一对，形成甲烷 \chem{CH_4}——天然气的主要成分，你家的燃气灶烧的就是它。碳尤其厉害：它还能和别的碳共享电子，连成链、连成环，搭出几百万种分子——那是整个第五篇（第 36 章起）要讲的故事，而这台大戏的每一块积木都是共价键。

\section{一对、两对、三对}

共享可以不止一对电子。两个原子之间可以共用两对甚至三对：

\begin{itemize}[itemsep=2pt]
\item \textbf{单键}：共享一对电子。两个氢原子之间、水分子里的 O—H 之间都是单键。
\item \textbf{双键}：共享两对电子。氧气分子里两个氧原子之间就是双键；你呼出的二氧化碳 \chem{CO_2} 里，碳和每个氧之间各是双键。
\item \textbf{三键}：共享三对电子。空气里最多的分子——氮气 \chem{N_2}——两个氮原子之间就是三键。
\end{itemize}

共享的电子对数目也叫\textbf{键级}：单键键级为 1，双键为 2，三键为 3。键级越高，两个核之间的电子云越“厚实”，原子被拴得越紧。这不是口号，是可以测量的规律，看一组真实数据：

\begin{table}[h]
\centering
\begin{tabular}{lccc}
\toprule
键 & 键级 & 键长（\ud{10^{-12}}{m}） & 键能（\ud{}{kJ\,mol^{-1}}） \\
\midrule
\chem{N-N}（单键） & 1 & 145 & 163 \\
\chem{N=N}（双键） & 2 & 125 & 418 \\
\chem{N \equiv N}（三键） & 3 & 110 & 945 \\
\chem{C-C}（单键） & 1 & 154 & 347 \\
\chem{C=C}（双键） & 2 & 134 & 614 \\
\chem{O=O}（氧气中） & 2 & 121 & 498 \\
\chem{H-H}（氢气中） & 1 & 74 & 432 \\
\bottomrule
\end{tabular}
\caption{几种共价键的键长与键能。键能指断开 1 摩尔这种键所需要吸收的能量（断键吸能，成键放能，见第 17 章）。数值是约数，同一类型的键在不同分子里会略有出入。}
\end{table}

趋势清清楚楚：键级越高，键越短、越强。氮—氮三键的键长比单键短了约四分之一，键能却是单键的近六倍。

\begin{qiaoliang}
键能到底有多大？拿氢分子算一笔账。断开 1 摩尔 H—H 键需要 \ud{432}{kJ} 的能量，1 摩尔包含 \ud{6.02\times 10^{23}}{} 个分子，所以拆开\textbf{一个}氢分子平均需要：
\[432\times 10^{3}\div 6.02\times 10^{23}\approx 7.2\times 10^{-19}\ \mathrm{J}\]
而室温下分子热运动的平均能量大约只有 $4\times 10^{-21}$ J——差了将近两百倍。这就是为什么室温下分子们撞来撞去，却撞不散共价键：碰撞递过来的那点钱，远远不够付“分手费”。反过来看，氮气三键的键能高达 \ud{945}{kJ\,mol^{-1}}，单个键要付的分手费约 $1.6\times 10^{-18}$ J，比氢分子还高出一倍多——所以氮气极其“懒”，空气里近八成的氮气天天从你身边流过，几乎什么都不参与。食品工厂正是利用这一点往薯片袋里充氮气当保护气：便宜、无毒、懒得反应。
\end{qiaoliang}

\section{拔河并不公平：电子云的偏移}

到现在为止，我们的例子像氢分子、氯分子，双方是同一个元素，力气完全一样，共享的电子对不偏不倚地待在中间。这种键叫\textbf{非极性共价键}。

但只要成键的是两种不同的原子，共享就几乎不可能绝对公平。第 11 章讲周期律时认识过\textbf{电负性}——一个原子在成键时把共享电子往自己这边拉的力气。氟的电负性最大，氧、氮、氯紧随其后，越往周期表左下走电负性越小。当两个电负性不同的原子共享电子对时，电子云会整体向电负性大的那一侧偏移：它在共享区域里“待的时间更长”。

结果是什么？电负性大的一端，电子云密度偏高，带上一点点负电（记作 $\delta^{-}$）；另一端电子被拉走一些，带上一点点正电（记作 $\delta^{+}$）。$\delta$ 是希腊字母，在这里表示“一丁点儿、不完整”的电荷——远不到离子那样整整一个单位的电荷。这种电子云分布不均匀的共价键叫\textbf{极性共价键}。

水分子是最好的例子。氧的电负性（3.4）明显大于氢（2.1），所以 O—H 键里电子云偏向氧：氧端带 $\delta^{-}$，氢端带 $\delta^{+}$。一个水分子一头微微偏负、两头微微偏正——请记住这个画面，它将在第 22、23 章掀起滔天巨浪：水的几乎所有怪脾气（沸点反常地高、冰能浮在水上、溶解能力极强）都从这“一点点偏”开始。

还有一个更深的认识：\textbf{非极性共价键、极性共价键、离子键不是三类东西，而是一条连续的滑梯}。电负性差为 0（同种原子），滑梯在最顶端，完全公平共享；差值增大，电子云越来越偏，就是极性键；差值大到一定程度（经验上大约超过 1.7），电子干脆被整个夺走，滑梯滑到底，变成第 18 章的离子键。食盐里钠和氯的电负性差约 2.2，是标准的离子键；水里的 O—H 差约 1.4，是极性相当明显的共价键；氯化氢 \chem{HCl} 差约 0.9，偏得中等。自然界没有在哪一格上划线说“这里以下才算共价键”——分类是我们为了方便贴的标签，滑梯本身是光滑的。

\begin{shiyan}
\textbf{会拐弯的水流。}材料：一把塑料梳子（或一个塑料尺）、你自己的头发、厨房水龙头。步骤：打开水龙头，把水流调到细细的一线（越细越好，不断流即可）；用梳子在干燥的头发上用力摩擦十几下，然后慢慢靠近细水流（不要碰到水）。观察点：水流会向梳子方向明显地弯曲，仿佛被一只看不见的手拉过去。多摩擦几次、换不同材质的梳子对比弯曲程度。解释：摩擦让梳子带上了静电；水分子一头偏正、一头偏负，靠近带电梳子时，水分子们纷纷转身、排队，异号的一端朝向梳子，整束水流就被吸弯了。如果水分子里的电子真是“平均分家”、不带任何偏性，水流根本不会理会梳子。安全提示：远离插座和电器，实验后擦干台面；仅此而已，这个实验很安全。
\end{shiyan}

\section{给分子画“结构图”：路易斯式登场}

道理讲了这么多，终于轮到符号上场了。1916 年，美国化学家路易斯发明了一种至今仍在全世界黑板上使用的画法：只画每个原子的\textbf{最外层电子}，用点（或叉）表示；两个原子之间共用的一对电子，后来简化成一条短线。这种图叫\textbf{路易斯结构式}。

画法只有三步：第一步，数出所有原子最外层电子的总数（这是“总预算”）；第二步，先用单键把原子连起来；第三步，把剩下的电子作为“孤对电子”补到各原子上，让每个原子（除氢只要 2 个外）周围尽量凑满 8 个电子——如果电子不够分，就把一对孤对电子“升级”为共享，形成双键或三键。

拿三个老朋友练手。氢气：总共 2 个电子，恰好一对，写成 H—H。水：氧带 6 个、两个氢各带 1 个，共 8 个电子；先用两条 O—H 单键用掉 4 个，剩 4 个作为两对孤对电子堆在氧上——氧周围 8 个，氢周围 2 个，预算刚好花完。二氧化碳：碳 4 个加两个氧各 6 个，共 16 个；先用两条单键用掉 4 个，剩下的电子怎么补都补不满碳的 8 个，于是把每个氧上的一对孤对电子拉过来共享，变成两条 C=O 双键，16 个电子刚好各就各位。

这套记账法还能算出一笔更细的账，叫\textbf{形式电荷}：把一个原子“名下”的电子数（孤对电子全归自己，共享电子对一人一半）和它\textbf{自己带来的}最外层电子数比一比，差额就是形式电荷。它回答的问题是：“在这张分家方案里，这个原子比它单干时多占了还是少占了？”以奇特的一氧化碳 \chem{CO} 为例：10 个最外层电子，唯一能让碳和氧都凑满 8 个的画法是三键加两对孤对电子。算账：氧在三键里分到 3 个、加自己的一对孤对 2 个，名下 5 个，比它带来的 6 个少 1 个，形式电荷 $+1$；碳名下 5 个，比带来的 4 个多 1 个，形式电荷 $-1$。结论有点反直觉：在一氧化碳里，居然是\textbf{电负性更大的氧带正的形式电荷}。形式电荷不是真实的电荷分布（真实电子云仍然偏向氧，CO 的极性很小但氧端仍略偏负），它只是记账工具——但这个工具很有用：合理的路易斯结构应当让形式电荷尽量小、尽量让负的形式电荷落在电负性大的原子上。用它来筛选“哪种画法更靠谱”，比瞎猜强得多。

还有一类分子，一张路易斯结构根本画不下。臭氧 \chem{O_3}：18 个最外层电子，画出来必然是一条双键加一条单键——可双键画在左边还是右边？两种画法都合法。而实验测得臭氧里两条 O—O 键\textbf{完全一样长}（约 \ud{128}{10^{-12}m}），介于单键和双键之间。出路是：真实分子不是两种画法中的任何一种，而是两者的“平均”——两对额外的电子\textbf{离域}在三个氧原子上，两条键各自相当于“一点五级”。用两张路易斯式加双头箭头并排表示，叫\textbf{共振结构}。请记牢：共振不是分子在两种结构之间来回切换、左右横跳——分子从头到尾只有一种、也是唯一的一种真实结构，只是我们的纸笔画不出来，只好画两张“近似照”让你脑补平均。纸笔的局限，不是分子的行为。

\section{科学家怎样知道：从立方原子到量子计算}

\begin{zhengju}
共价键的概念诞生时，其实是一段“没有理论、只有账本”的岁月。1916 年，路易斯提出：原子可以靠共享电子对达到外层 8 电子的排布。他当时连原子的内部结构都不清楚（电子云、轨道都是十年后的概念），他的原子模型甚至是静止的“立方原子”——八个电子蹲在立方体的八个角上，两个立方体共用一条边就是共享一对电子。这个模型今天看近乎可笑，但它第一次统一解释了一大堆事实：为什么氢、氧、氮都是双原子分子，为什么碳能连成长链，为什么分子的组成总遵守简单的整数比。好的模型不必一开始就正确，它先要\textbf{好用}。

但路易斯回答不了“为什么共享能让能量降低”。在他的图景里，两个带负电的电子凑在一起只会互相排斥，共享一对电子凭什么让两个原子粘住？反对者完全有理由不买账。当时流行的另一种理论把一切都归结为静电（离子理论），它解释食盐很成功，可面对两个一模一样的氢原子就彻底哑火——两边电荷相同，静电从哪来？

答案要等量子力学。1927 年，两位物理学家海特勒和伦敦把刚诞生两年的薛定谔方程用到氢分子上：他们写下两个氢原子的电子波函数，计算体系能量随核间距的变化。计算结果真的出现了一条带谷底的曲线，谷底位置（约 \ud{74}{10^{-12}m}）和深度（约 \ud{432}{kJ\,mol^{-1}}）与实验测量惊人吻合。更关键的是，能量下降的主要来源是一个经典物理里根本不存在的“交换项”：当两团电子云重叠时，电子无法分辨“自己属于哪个核”，两个电子交换位置后的状态必须叠加进来一起算——正是这个纯粹量子的效应把能量压了下去。共价键不是小棍，也不是简单的静电，它是量子力学写在分子里的条款。这是科学史上第一次，化学键从“经验口诀”变成了“可以从方程里算出来的东西”。

接力棒的下一程由鲍林完成。他把量子力学的语言翻译回化学家能用的工具：杂化、共振、以及 1932 年提出的电负性标度——就是我们前面用来判断“电子云往哪边偏”的那把尺子。他的《化学键的本质》影响了几代化学家。不过故事也有教训：鲍林晚年在维生素 C 等问题上过度自信、越过了证据边界，这提醒我们——再伟大的工具箱，也要一件一件地接受证据检验。
\end{zhengju}

\section{这套工具何时会失灵}

路易斯结构是初中生能掌握的最强工具之一，但它是\textbf{账本，不是照片}。它的适用边界，你迟早会踩到：

\begin{itemize}[itemsep=2pt]
\item \textbf{奇数电子的分子画不出完美结构}。一氧化氮 \chem{NO} 只有 11 个最外层电子，怎么画都剩一个“单身”电子。这不是你画错了，是分子真的带着一个未成对电子——这种分子（自由基）格外活泼，你体内每时每刻都在产生和清除它们（第 53 章还会再见）。
\item \textbf{有的原子宁愿不凑满 8 个}。三氟化硼 \chem{BF_3} 里，硼周围只有 6 个电子，照样稳定存在。缺电子不等于不稳定——能量说了算，口诀说了不算。
\item \textbf{有的原子超过 8 个也没事}。六氟化硫 \chem{SF_6} 里硫周围摊着 12 个电子。八隅体规则从第三周期往下越来越靠不住。
\item \textbf{共振分子画不出真身}。臭氧、硝酸根、苯环，真实结构都是几张路易斯式的“平均”，任何单独一张都在撒谎（苯环的故事在第 41 章）。
\item \textbf{金属完全不吃这一套}。铜、铁里的原子不靠一对一对地分电子，而是把电子撒进一片“公共海洋”——那是第 21 章的金属键。
\end{itemize}

还有最要紧的一条边界，也是本章标题的另一半：共享的电子\textbf{不是两颗固定在原子之间的小珠子}。路易斯式上的那条短线、我们图上的那对点，都只是记账符号。真实的图景是第 9 章的电子云：电子以概率的形式弥漫在两个原子核周围，问“这一刻这个电子在哪”是没有确定答案的。共价键的本质，是电子云重叠之后整个体系能量降低这件事本身——珠子、小棍、短线，全是我们为了方便发明的人为表示法。

\begin{wuqu}
“共价键就是两个原子一人出一个电子，公平分家、谁也不吃亏。”——两个错误打包在一起。其一，共享不等于平均：只要两种原子的电负性不同，电子云就偏向力气大的一边。水分子里的氧把共享电子拉得离自己更近，于是氧端偏负、氢端偏正——带电梳子能吸弯水流，就是“分家不公平”的直接证据（见本章实验）。其二，“一人出一个电子”也不是定律：在配位键里，成键的一对电子可以由同一方全部提供，比如一氧化碳里就可以看作氧多“出了”一对（这正是它形式电荷反常的原因），铵离子 \chem{NH_4^{+}} 里的第四条 N—H 键也是氮单独提供的电子对——成键之后，这条键和其他三条完全无法区分。电子只服从能量规律，从不遵守“公平”。
\end{wuqu}

\section{回到那只蓝色钢瓶}

现在重新解释开场的一切。

钢瓶里的氧气：两个氧原子，谁也抢不走谁的电子，于是共享两对，结成双键 \chem{O=O}，键能约 \ud{498}{kJ\,mol^{-1}}——足够牢，所以氧气能压缩存放几年；又不够牢，所以在你的细胞里，它会被一台精密的分子机器一点点拆开、驱动整个能量供应（第 53 章）。

那杯水：每个水分子内部，氧和氢靠极性共价键拴在一起，电子云偏向氧——所以水分子一头偏负、一头偏正，能被带电梳子拉动。烧水时你并没有拆开水分子（那要上千度），只是给分子们足够能量挣脱彼此之间的松散吸引（第 22、23 章），所以水烧开了还是水。

塑料笔杆和头发：它们是碳、氢、氧、氮等元素用共价键连成的长链分子，一条链上可能有成千上万个共价键。拉扯时链子会先滑动、绷直，所以塑料先变形才断裂——长链分子世界的故事在第 43 章。至于橡皮筋能拉长五倍弹回来，靠的也是共价键搭的碳链骨架。

你看，从钢瓶到发丝，撑起这一切的不过是同一条规则：电子云重叠，体系能量降低，原子便结伴。

\section{结伴之后，还要排队}

这一章解决了“原子怎样共享电子”，但立刻冒出下一个问题：共享电子对和孤对电子都带负电，它们在一个分子里会怎样安排自己的位置？

想想水：氧周围有两条 O—H 键和两对孤对电子，四团电子云互相排斥，绝不会排成一条直线——所以水分子是弯的。二氧化碳里的两条双键则尽量躲到相反方向，所以 \chem{CO_2} 是笔直的。甲烷的四条键彼此撑到最开，指向正四面体的四个角。分子的形状不是随机长成的，它是电子对之间排斥的必然结果——而形状又直接决定分子能做什么：水的弯曲让它的正负两端“露在外面”， \chem{CO_2} 的笔直让两条极性键的偏移恰好抵消；你的嗅觉、酶的催化、DNA 的双螺旋，全都建立在分子形状的精密配合上。这就是下一章的主题：\textbf{分子的形状决定它会做什么}。

\begin{tiaozhan}
为什么共价键总是“成对”共享，而不是三个、五个电子一起共享？答案在电子的一种量子属性——自旋。电子像一枚只能取两种朝向的小磁针（向上或向下），而泡利不相容原理规定：同一个“房间”（轨道）里最多住两个电子，且必须自旋相反。海特勒—伦敦的计算显示，只有当两个电子自旋相反时，交换效应才把能量压下去形成谷底；自旋相同时能量反而升高，原子互相推开。所以成键的电子天然是一对一对的。再进一步：鲍林之外的另一套理论（分子轨道理论，由马利肯等人发展）干脆不承认电子“属于哪条键”，而是把整个分子当成一个整体来安放电子——它能轻松解释一些路易斯理论说不通的事实，比如氧气分子里其实住着两个未成对电子（液氧能被磁铁吸住！），而路易斯式画出来的 \chem{O=O} 全体成对。两套理论谁对？都对，都是近似，适用场合不同——这正是科学模型的常态。跳过本节不影响主线。
\end{tiaozhan}

\section*{练一练}

\begin{sikaoti}
\item 画出氨分子 \chem{NH_3} 的路易斯结构（氮带 5 个最外层电子，每个氢带 1 个）。数一数：氮周围有几对成键电子、几对孤对电子？总电子预算是多少，花完了吗？
\item 用能量观点解释：为什么说“断开共价键需要吸收能量”而不是“断键释放能量”？请画一条两个原子的势能曲线（纵轴为能量，横轴为核间距），在图上标出键长和键能的位置。
\item 查一查电负性数值，把下列键按“电子云偏移程度”从小到大排队：C—H、O—H、H—H、Na—Cl、H—Cl。哪一对已经滑出共价键的范围，更像离子键？
\item 有人说“双键就是两条单键，所以强度应该是单键的两倍”。用本章表格里 \chem{C-C}（347）和 \chem{C=C}（614）的数据检验这个说法。如果它不对，猜猜第二条键为什么“不如第一条结实”（提示：两对电子塞在同一对原子之间，会互相排斥）。
\item 硝酸根 \chem{NO_3^{-}} 可以画出三张等价的共振结构（双键轮流落在三个氧上）。由此预测：三个 N—O 键的键长会不会相同？三个氧原子分摊的负电荷各是多少？
\item 液氧可以被强磁铁吸住，说明氧气分子里有未成对电子；但路易斯结构画出的 \chem{O=O} 里所有电子都成对。这个矛盾说明路易斯结构这个工具怎么了？它应该被扔掉吗，还是应该被重新定位？
\end{sikaoti}
